% report/sections/profiling.tex
% Author: Bhanuja Karumuru  |  Milestones: M9, M12
% Hardware design, NSight profiling methodology, SM occupancy analysis, figures.

\section{Hardware Design and Profiling}
\label{sec:profiling}

\subsection{Tiling Strategy and Shared Memory Layout}
\label{sec:tiling}

The target hardware is an NVIDIA A100-SXM4-40GB (NYU Greene cluster),
which provides 164\,KB of configurable shared memory per SM across 108 SMs,
with HBM bandwidth of 1{,}555\,GB/s.

\paragraph{Shared memory budget.}
Each CUDA block handles one $(\text{batch}, \text{head})$ pair
(grid = $(B, H)$, block = $(T, 1)$).
The block holds three shared-memory arrays --- $s_Q$, $s_K$, $s_V$ ---
each of shape $T \times (d_{\text{head}} + 1)$ in fp32.
(The output accumulator $s_O$ is kept in thread-private registers to reduce
shared memory pressure.)
Total shared memory per block:
\begin{equation}
  \text{shmem} = 3 \times T \times (d_{\text{head}} + 1) \times 4~\text{bytes.}
\end{equation}
For the profiling benchmark ($d_{\text{head}} = 64$, padded stride = 65):

\begin{table}[h]
\centering
\caption{Tile-size shared-memory budget and theoretical maximum concurrent blocks per SM (A100, $d_{\text{head}}=64$).}
\label{tab:shmem_budget}
\begin{tabular}{rrrrc}
\hline
$T$ & shmem/block (bytes) & shmem/block (KB) & max blocks/SM & \\
\hline
 16 & 12,480 & 12.2 & 13 & \\
 32 & 24,960 & 24.4 &  6 & \\
 \textbf{64} & \textbf{49,920} & \textbf{48.8} & \textbf{3} & \textbf{selected} \\
128 & 99,840 & 97.5 &  1 & \\
\hline
\end{tabular}
\end{table}

$T = 64$ was selected: it uses 49\,KB (30\% of the 164\,KB budget),
supports 3 resident blocks per SM to hide latency,
and avoids the register-pressure spike observed at $T = 128$
which constrains occupancy to 1 block/SM.
The M9 occupancy sweep (\texttt{results/occupancy\_sweep.csv}) confirms
this empirically via NSight Compute.

\paragraph{Bank-conflict elimination.}
\label{par:bank_conflict}
The A100 shared memory has 32 banks, each 4 bytes wide.
A na\"ive array declared as \texttt{float arr[T][64]} has inner stride 64,
which is an exact multiple of 32.
For a fixed column $c$, all 32 threads in a warp map to bank
$(64i + c) \bmod 32 = c \bmod 32$ --- a 32-way bank conflict.

The fix is to pad the inner dimension by 1:
\begin{lstlisting}[language=C]
__shared__ float sQ[TILE_SIZE][HEAD_DIM + 1];  // [64][65]
__shared__ float sK[TILE_SIZE][HEAD_DIM + 1];
__shared__ float sV[TILE_SIZE][HEAD_DIM + 1];
\end{lstlisting}
With stride 65 ($65 \bmod 32 = 1$), thread $i$ accessing column $c$ maps to
bank $(65i + c) \bmod 32 = (i + c) \bmod 32$.
The 32 warp threads access banks
$c\bmod 32,\; (c{+}1)\bmod 32,\; \ldots,\; (c{+}31)\bmod 32$ ---
all distinct, eliminating conflicts entirely.
The padding overhead is $3 \times T \times 4 = 768$ bytes at $T = 64$,
less than 2\% of the shared memory allocation.

\paragraph{Thread organisation and coalesced HBM access.}
Thread \texttt{tid} (0--63) owns token row \texttt{tid} within the current tile.
Loading $s_Q[\texttt{tid}][\cdot]$ assigns consecutive threads to consecutive
token rows in HBM, so the 32-thread warp reads a 32-row contiguous block
--- one 128-byte cache transaction per 32 float elements,
with zero wasted bandwidth.
The same pattern holds for $s_K$, $s_V$, and the output write.
Zero warp divergence occurs in the compute phase: every thread executes the
same dot-product loop body with no data-dependent branches.

\paragraph{\texttt{\_\_syncthreads()} placement.}
Two barriers are inserted per K/V-tile iteration:
(1) immediately after loading $s_K$ and $s_V$, before any score computation;
(2) after all threads finish reading $s_K$/$s_V$, before the next iteration
    overwrites them.
One additional barrier separates each Q-tile's output write from the next
Q-tile's load.
Omitting barrier (1) allows a thread that finishes its load early to read
stale data written by slower threads; omitting barrier (2) allows the next
iteration to corrupt shared memory that is still being read.

\subsection{Profiling Methodology}
\label{sec:profiling_method}

All profiling experiments mirror the baseline protocol (Rithwik, M3) exactly.

\begin{itemize}
  \item \textbf{Hardware:} NVIDIA A100-SXM4-40GB, NYU Greene cluster.
  \item \textbf{Configuration:} embed\_dim = 512 (\texttt{config.EMBED\_DIM\_BENCH}),
        $H = 8$ (\texttt{config.N\_HEADS\_BENCH}),
        $d_{\text{head}} = 64$ (\texttt{config.D\_HEAD}),
        batch = 1 (\texttt{config.BATCH\_BENCH}).
  \item \textbf{Sequence lengths:} $N \in \{64, 128, 256, 512, 1024\}$.
  \item \textbf{Warmup:} 100 kernel invocations (GPU reaches steady thermal state).
  \item \textbf{Timed iterations:} 500, bracketed by \texttt{torch.cuda.Event}
        start/stop; average per-iteration time reported.
  \item \textbf{NSight Systems} (\texttt{nsys profile --trace=cuda,nvtx}):
        captures the full kernel launch timeline, confirming single vs.\ dual
        kernel per forward pass.
  \item \textbf{NSight Compute} (\texttt{ncu --set full}):
        exact HBM bytes read/written (metric \texttt{l1tex\_\_t\_bytes}),
        SM occupancy (\texttt{sm\_\_warps\_active.avg.pct\_of\_peak\_sustained\_active}),
        and memory throughput. Run per sequence length with 10 warm + 10 timed
        iterations to avoid trace-file bloat.
\end{itemize}

The fused benchmark script (\texttt{profiling/fused\_bench.py}) uses identical
warmup counts, timing protocol, and CSV schema as \texttt{profiling/baseline\_bench.py}
so that \texttt{results/merge\_comparison.py} can directly compute
speedup = baseline\_\textmu s / fused\_\textmu s per row.

\subsection{Results}
\label{sec:profiling_results}

\paragraph{Kernel timeline.}
Figure~\ref{fig:nsight_timeline} shows side-by-side NSight Systems kernel timelines.
The baseline requires two separate CUDA kernel launches per forward pass
(QKV projection followed by SDPA), with HBM round-trips for $Q$, $K$, $V$
in between.
The fused kernel replaces both with a single launch.

\begin{figure}[h]
\centering
\includegraphics[width=\linewidth]{results/figures/nsight_timeline.png}
\caption{NSight Systems kernel timeline. \emph{Top}: baseline two-kernel pipeline
(QKV GEMM then SDPA); \emph{bottom}: fused single-kernel.
The gap between kernels in the baseline represents the HBM traffic carrying
$Q$, $K$, $V$ between them.}
\label{fig:nsight_timeline}
\end{figure}

\paragraph{Baseline timing (Rithwik, M3).}
The measured per-iteration wall times on the A100 are:

\begin{table}[h]
\centering
\caption{Baseline (unfused) per-iteration wall time, A100-SXM4-40GB.}
\label{tab:baseline_timing}
\begin{tabular}{rr}
\hline
$N$ & per-iteration (\textmu s) \\
\hline
  64 & 302.8 \\
 128 & 305.5 \\
 256 & 301.1 \\
 512 & 270.3 \\
1024 & 379.6 \\
\hline
\end{tabular}
\end{table}

(Source: \texttt{results/baseline\_profiling.csv}.)

\paragraph{Fused kernel timing and speedup.}
Table~\ref{tab:comparison} will be populated after M9 profiling runs
and merged via \texttt{results/merge\_comparison.py}.
Predicted speedup range: 1.3$\times$--2.5$\times$ based on the HBM traffic
reduction estimated in \texttt{kernel/DESIGN.md}.

\begin{table}[h]
\centering
\caption{Baseline vs.\ fused: wall time and kernel count, all sequence lengths.}
\label{tab:comparison}
% Generated by: python results/merge_comparison.py
% Source: results/comparison_table.csv
\input{results/comparison_table_tex.tex}
\end{table}

\paragraph{HBM bandwidth.}
Figure~\ref{fig:hbm} shows HBM read traffic for baseline and fused kernels.

\begin{figure}[h]
\centering
\includegraphics[width=0.85\linewidth]{results/figures/hbm_bandwidth.png}
\caption{HBM read traffic (bytes) vs.\ sequence length, baseline (unfused) vs.\ fused.
The fused kernel eliminates intermediate $Q$/$K$/$V$ HBM writes and re-reads.}
\label{fig:hbm}
\end{figure}

\paragraph{Wall-time speedup.}
Figure~\ref{fig:speedup} plots speedup as a function of $N$.
Speedup grows with sequence length because the eliminated HBM traffic
scales as $\mathcal{O}(N)$ while compute scales as $\mathcal{O}(N^2)$;
at longer sequences the kernel spends a larger fraction of time on memory
traffic that the fused design eliminates.

\begin{figure}[h]
\centering
\includegraphics[width=0.85\linewidth]{results/figures/speedup.png}
\caption{Per-iteration wall-time speedup (baseline / fused) vs.\ sequence length.}
\label{fig:speedup}
\end{figure}

\paragraph{SM occupancy and tile-size sweep.}
Figure~\ref{fig:occupancy} shows the occupancy--throughput trade-off across
tile sizes $T \in \{16, 32, 64, 128\}$ at $N = 512$.
Full data is in \texttt{results/occupancy\_sweep.csv}.

\begin{figure}[h]
\centering
\includegraphics[width=0.85\linewidth]{results/figures/occupancy_vs_tile.png}
\caption{SM occupancy (\%) and per-iteration wall time (ms) vs.\ tile size $T$,
$N = 512$. $T = 64$ achieves the best throughput--occupancy balance.}
\label{fig:occupancy}
\end{figure}
